\section{Discussion}
\label{sec:discussion}

\subsection{What the unification buys}

Three things. First, economy: a single theorem and a single rank
condition replace five named theorems and five rank arguments, and a
result proved once at the right generality is easier to extend than
five parallel proofs. A new domain that fits the abstract process of
\cref{sec:framework} inherits consistency, identifiability, and the
singleton remedy without a fresh derivation; the only domain-specific
work is to identify the latent variable, the report, the
coarsening-sufficient statistic, and the singleton, which is the
content of \cref{tab:css,tab:singletons}.

Second, a shared diagnostic moral. The single sentence that
\cref{thm:general-consistency} licenses, marginal-fit agreement is not
evidence of unbiasedness, is exactly the warning each sibling issues in
its own terms. Stating it once makes clear that it is a property of
coarsened-data MLE in general, not a quirk of dropout or of noised
releases or of imperfect codes. The practitioner check ``does my fitted
model reproduce the observed marginal'' is structurally uninformative
about the latent parameter in every one of these settings, and the only
cures are the singleton device or external information.

Third, transfer of remedies. Because the singletons are one mechanism
(\cref{prop:singleton}), a technique developed for one field's singleton
can in principle be ported to another's: the spike-in-fit bias analysis
of \citet{towell2026scrnacoarsening} and the gold-set sample-complexity
analysis of \citet{towell2026weaksupcoarsening} are asking the same
question, how much singleton information is needed, and their answers
are comparable through the rank deficit.

\subsection{The boundary of the principle}

A unification is most useful when its boundary is located exactly, not
asserted. The consistency identity has two places where the reach is
sharper than a blanket exact-for-all-domains claim, and both are
characterized rather than left open.

\paragraph{Differential privacy is a location family, not an
exponential family.} The other four consistency theorems are the
exponential-family mean-value identity (regime A). Release consistency
is a first-moment identity for a continuous convolution (regime B),
proved through the location-family score. The two are instances of one
stationarity principle, but regime (B) is not a corollary of regime
(A). Within regime (B) the boundary is exact: the location MLE equals
the arithmetic sample mean for every sample if and only if the kernel
is Gaussian (\cref{rem:loc-sketch}), the Cauchy-equation characterization
whose classical face is Gauss's characterization of the normal law.
Release consistency lives at the single release ($n = 1$), where the
identity is exact for any symmetric unimodal kernel; for $n \ge 3$ and a
non-Gaussian kernel the sample-mean form picks up a genuine
$O_p(n^{-1/2})$ remainder, and the Gaussian case is already covered by
regime (A) because the Gaussian location family is an exponential
family. There is nothing left open here: the finite-sample failure is
exactly the failure of the score to be linear.

\paragraph{Weak supervision reduces exactly only under a richer
parametrization.} Agreement consistency is exact when the agreement
indicators are sufficient statistics, but the naive-Bayes label model
that data programming actually uses does not make them sufficient, and
there the identity is asymptotic (\cref{cor:weaksup}). The unification
holds in both cases; the exact statement is for a parametrization
practitioners do not use, and for the one they do use the result is a
limit theorem.

A third, milder boundary: the spatial spot-level identity is exact per
coordinate but its vector form leans on the joint rank condition
(\cref{cor:spatial}), so consistency and identifiability are not as
cleanly separable there as in the scalar domains.

\subsection{Open per-domain problems}

The synthesis inherits the siblings' open problems and locates them in
the common structure.
\begin{itemize}
\item \textbf{Weak supervision, the rank-deficit factor.} The gold-set
  sample complexity is settled at $\Theta(r/\mathrm{gap}^2)$ for total
  (L2) model recovery, with matching upper and minimax lower bounds, in
  the parameter-degeneracy dimension $r$ and the accuracy margin; the
  rate is loss-dependent (it falls to $\Theta(\log r/\mathrm{gap}^2)$
  under a per-direction criterion). What remains open is estimating the
  degeneracy dimension $r$ from data and the multiclass extension
  \citep{towell2026weaksupcoarsening}.
\item \textbf{Phenotyping, real data and chart-review error.} The
  chart-reviewed state is treated as a true singleton; adjudication is
  itself imperfect, which is a relaxation of the singleton assumption,
  and a real-data validation remains to be done
  \citep{towell2026phenotypecoarsening}.
\item \textbf{Differential privacy, finite-sample multi-release
  corrections.} The single-release identity is exact and the
  Gaussian-iff characterization of the sample-mean form is settled
  (\cref{rem:loc-sketch}); what remains optional, not open, is
  tabulating the leading constant $V = \operatorname{Var}(\psi(Z)/J - Z)$
  of the $O_p(n^{-1/2})$ remainder for the standard release kernels if a
  finite-sample correction over several releases is wanted.
\item \textbf{C2 violations, quantitatively.} The sensitivity bands of
  \citet{towell2026mdrelax} are worked out for reliability; porting the
  bounded-C2-violation analysis to the other domains is open and would
  give each a robustness counterpart to its consistency theorem.
\end{itemize}

\subsection{Relation to the coarsening-at-random literature}

The framework is a specialization of coarsening at random
\citep{heitjan1991ignorability,gill1997coarsening} and of the
missing-data ignorability tradition \citep{little2002statistical,tsiatis2006semiparametric}.
What is new here is not the CAR conditions, which are classical, but the
observation that one consistency identity and one rank condition recur
across six concrete fields as instances of those conditions, and the
explicit reduction of the five named theorems to a single statement
with a clean account of the two regimes and the two seams. The
exponential-family mean-value identity that powers regime (A) is itself
textbook \citep{vandervaart1998asymptotic}; the contribution is to
recognize that the five named theorems are that identity wearing five
costumes.
